% Part 6 导言：应用领域
\vspace{1em}
\begin{center}
\textcolor{techgray}{\small 本部分展示3DGS在六个垂直领域的深度应用——从数字人到安全隐私，每个领域都有独特的技术挑战和解决方案。}
\end{center}
\vspace{0.5em}

\noindent\textbf{本部分包含：}
\begin{itemize}[leftmargin=2em,itemsep=1pt]
\item[\textcolor{techgreen}{\textbullet}] \textbf{第16章}（人体化身）— 人脸/身体建模、表情驱动、语音驱动动画
\item[\textcolor{techgreen}{\textbullet}] \textbf{第17章}（机器人）— 导航规划、机械臂操控、具身感知、仿真到真实
\item[\textcolor{techgreen}{\textbullet}] \textbf{第18章}（自动驾驶）— 街景重建、仿真、动态场景、传感器融合
\item[\textcolor{techgreen}{\textbullet}] \textbf{第19章}（医学影像）— 内窥镜重建、CT/MRI、显微成像、临床应用
\item[\textcolor{techgreen}{\textbullet}] \textbf{第20章}（遥感测绘）— 卫星重建、航空影像、文化遗产、测绘
\item[\textcolor{techgreen}{\textbullet}] \textbf{第21章}（安全隐私）— 数字水印、对抗攻防、隐写
\end{itemize}

\noindent\textbf{阅读建议：}各章完全独立，可直接跳到感兴趣的领域。第16章需要第5章（材质）基础；第17章需要第15章（SLAM）基础；第18章需要第6章（动态）和第14章（大规模）基础。

\noindent\textbf{与其他部分的联系：}六个应用章节是前五部分技术成果的综合体现——每章都会回溯引用前面章节的核心方法，读者可在阅读应用章节时反向追踪感兴趣的技术细节。

\vspace{1em}
